\documentclass{article}
\usepackage{graphicx} % Required for inserting images
\usepackage{amsmath}
\usepackage{svg}
\usepackage{listings}
\usepackage{tcolorbox}
\tcbset{
    cppstyle/.style={
        colback=gray!10,    % Background color
        colframe=gray!50,   % Border color
        fonttitle=\bfseries,% Title font
        coltitle=black,     % Title color
        sharp corners,      % Square corners
        boxrule=0.5mm,      % Thickness of the border
        left=2mm,           % Left padding
        right=2mm,          % Right padding
        top=1mm,            % Top padding
        bottom=1mm,         % Bottom padding
        boxsep=2mm,         % Space around the box
    }
}

\usepackage{xcolor} % Optional, for custom colors

\title{trial}
\author{Jay Bhambure}
\date{June 2024}

\begin{document}

\section{Main Ideas:}
The main objective of the vischydro.cpp code is to simulate the evolution of a hydrodynamic system using ideal hydrodynamics. The code uses PETSc (Portable, Extensible Toolkit for Scientific Computation) to handle parallel computations and solve the hydrodynamic equations over time.

Key Components and Their Roles \\
\subsection{Equation of State (EOS) Class:
}
The \textbf{\textit{EOS}} class provides methods to calculate various thermodynamic properties such as pressure, temperature, and speed of sound squared based on the energy density and baryon density.
\subsection{2. VischydroNode Structure:
}
The \textit{\textbf{VischydroNode}} structure represents a cell in the hydrodynamic simulation. It contains properties like energy, momentum, pressure, and velocity.
\subsection{3. Hydrodynamic Functions:
}
\textit{\textbf{FillVischydroNode}}: Initializes a \textit{\textbf{VischydroNode}} with values from the \textit{\textbf{EOS}}.
\textbf{\textit{idealHydroCellIFunction}} and \textbf{\textit{idealHydroCellIFunctionDerivative}}: Functions used to check consistency and calculate derivatives for the Newton-Raphson method.
\textbf{\textit{idealHydroCellSolve}}: Uses the Newton-Raphson method to solve for the energy density in a hydrodynamic cell.
\subsection{4. Vischydro Class:
}
The Vischydro class manages the hydrodynamic simulation. It sets up the computational domain, initializes the solution vector, and handles time-stepping using PETSc's TS (time-stepping) routines.
It reads initial conditions from an HDF5 file, sets up the time-stepping parameters, and advances the solution in time.
\subsection{5. PETSc Integration:
}
The code uses PETSc to handle parallel computations, solve differential equations, and manage data.
Functions like \textbf{\textit{EulerRHSFunction}}, \textbf{\textit{LHSIFunction}}, and \textbf{\textit{LHSIJacobian}} define the right-hand side, implicit function, and Jacobian for the time-stepping solver.
\subsection{6. Main Function:
}
The main function initializes PETSc, reads input parameters from a JSON file, sets up the Vischydro object, and advances the solution in time using \textbf{\textit{TSSolve}}.
It also handles writing the solution to an HDF5 file at each time step.
\subsection{Inputs and Outputs
}\textbf{Inputs}: JSON configuration file (inputs.json or specified via command line) containing parameters like grid size, initial time, CFL condition, and final time.
HDF5 file containing initial conditions for the hydrodynamic variables.\\
\textbf{Outputs:} HDF5 file containing the solution at each time step and the final solution.
Summary\\

The vischydro.cpp code simulates the evolution of a hydrodynamic system using ideal hydrodynamics. It uses PETSc to handle parallel computations and solve the hydrodynamic equations over time. The main steps include reading input parameters, initializing the computational domain and solution vectors, advancing the solution in time using PETSc's time-stepping routines, and writing the solution to an HDF5 file.

\lstinline|int x = 5;|
\end{document}
